\documentclass[a4paper,12pt]{article}
\usepackage[utf8]{inputenc}
\usepackage[T1]{fontenc}
\usepackage[ngerman]{babel}
\usepackage{amsmath}
\usepackage{amssymb}
\usepackage{enumitem}
\usepackage{geometry}
\geometry{a4paper, margin=2.5cm}

\title{Einsendeaufgaben -- Aufgabe 2.1\\[0.5em]
\large 61211 Analysis}
\author{}
\date{}

\begin{document}
\maketitle

\section*{Aufgabe 1}

\begin{enumerate}[label=\alph*)]

    \item
    Da für alle $x_0 \in \mathbb{R}$ $1 + x_0^2 \neq 0$ gilt, können wir Satz 2.3.8(iv) anwenden.
    $f$ ist also stetig.

    \item
    $f$ ist nicht stetig. Wir zeigen dies, indem wir zeigen, dass $f$ nicht das Folgenkriterium erfüllt.

    Es gilt $f(0,0) = 0$.

    Wir definieren eine Folge $(x_k)$ mit $x_k := \left(\frac{1}{k}, \frac{1}{k}\right)$, die gegen 0 in $(\mathbb{R}^2, \|\cdot\|_2)$ konvergiert.

    Es gilt allerdings nicht, dass $f(x_k) \to f(0,0) = 0$ in $(\mathbb{R}, |\cdot|)$, denn
    \[
        \lim_{k \to \infty} f\!\left(\tfrac{1}{k}, \tfrac{1}{k}\right)
        = \lim_{k \to \infty} \frac{\left(\tfrac{1}{k}\right)^2 + \left(\tfrac{1}{k}\right)^2}{\left(\tfrac{1}{k}\right)^2}
        = \lim_{k \to \infty} 2 = 2 \neq 0.
    \]

    Demnach ist $f$ in 0 nicht stetig.

    \item
    Nach der lokalen Eigenschaft der Stetigkeit (Satz 2.3.4(ii)) und Satz 2.3.8(iv) ist $f$ in jedem
    Punkt $a \in \mathbb{R}^2$ mit $a \neq (0,0)$ stetig.

    Wir untersuchen noch den Fall $a = 0$. Wir werden das $\varepsilon$-$\delta$-Kriterium anwenden.

    Sei $\varepsilon > 0$ beliebig. Wir wählen $\delta := \dfrac{\varepsilon}{2}$.

    Sei $x \in \mathbb{R}^2$ mit $\|x\|_2 \leq \delta$ beliebig. Dann gilt:
    \begin{align*}
        \|f(x) - f(0)\| &= \|f(x)\| \\
        &= \left|\frac{x_1^2 + x_1 x_2 + x_2^2}{\sqrt{x_1^2 + x_2^2}}\right| \\
        &\leq \left|\frac{x_1^2 + x_2^2}{\|x\|_2}\right| + \left|\frac{x_1 x_2}{\|x\|_2}\right|
    \end{align*}

    Es gilt $|x_1 x_2| \leq x_1^2 + x_2^2$, da:

    \textit{Fall 1:} $x_1 x_2 < 0$:
    \[
        x_1^2 + x_1 x_2 + x_2^2 \geq x_1^2 + 2x_1 x_2 + x_2^2 = (x_1 + x_2)^2 \geq 0
    \]

    \textit{Fall 2:} $x_1 x_2 \geq 0$:
    \[
        x_1^2 - x_1 x_2 + x_2^2 \geq x_1^2 - 2x_1 x_2 + x_2^2 = (x_1 - x_2)^2 \geq 0
    \]

    Daraus können wir schließen:

    \begin{align*}
        \|f(x) - f(0)\| \\
        &\leq \left|\frac{x_1^2 + x_2^2}{\|x\|_2}\right| + \left|\frac{x_1 x_2}{\|x\|_2}\right| \\
        &\leq \left|\frac{x_1^2 + x_2^2}{\|x\|_2}\right| + \left|\frac{x_1^2 + x_2^2}{\|x\|_2}\right| \\
        &= 2\left|\frac{\|x\|_2^2}{\|x\|_2}\right| < 2 \cdot \frac{\varepsilon}{2} = \varepsilon
    \end{align*}

\end{enumerate}

\end{document}
